%%%%%%%%%%%%%%%%%%%%%%%%%%%%%%%%%%%%%%%%%%%%%%%%%%%%%%%%%%%%
% Rating Methodologies
%%%%%%%%%%%%%%%%%%%%%%%%%%%%%%%%%%%%%%%%%%%%%%%%%%%%%%%%%%%%

\section{Rating Agency Methodologies}

Credit Rating Agencies (CRAs) play a crucial role in the securitisation process. They evaluate the credit quality of loan pools and assign ratings that help investors understand the level of credit risk. A higher rating indicates a lower probability of default, while a lower rating indicates higher investment risk.

\vspace{0.4cm}

\subsection*{Major Rating Agencies}

\renewcommand{\arraystretch}{1.5}

\begin{tabularx}{\textwidth}{|p{3.5cm}|X|}
	\hline
	\rowcolor{TableHeader}
	\textbf{Rating Agency} & \textbf{Methodology} \\
	\hline
	
	\textbf{CRISIL}
	&
	Evaluates asset quality, historical default trends, recovery rates, portfolio diversification, and credit enhancement before assigning ratings.
	\\
	\hline
	
	\textbf{Moody's}
	&
	Uses the Idealised Default Rate (IDR) model to estimate expected default probability and expected loss over different rating categories.
	\\
	\hline
	
	\textbf{S\&P Global}
	&
	Uses the LEVELS methodology to analyse portfolio quality, cash flow strength, macroeconomic risk, legal structure, and servicing quality.
	\\
	\hline
	
\end{tabularx}

\vspace{0.6cm}

\subsection*{Typical Rating Process}

\begin{center}
	
	\fbox{
		\begin{minipage}{0.90\textwidth}
			
			\centering
			
			Loan Pool
			
			$\downarrow$
			
			Data Quality Review
			
			$\downarrow$
			
			Default Probability Analysis
			
			$\downarrow$
			
			Cash Flow Modelling
			
			$\downarrow$
			
			Stress Testing
			
			$\downarrow$
			
			Credit Enhancement Assessment
			
			$\downarrow$
			
			Final Credit Rating
			
		\end{minipage}
	}
	
\end{center}

\vspace{0.6cm}

\subsection*{Comparison of Rating Methodologies}

\renewcommand{\arraystretch}{1.4}

\begin{tabularx}{\textwidth}{|p{4cm}|X|X|X|}
	\hline
	\rowcolor{TableHeader}
	\textbf{Criteria} & \textbf{CRISIL} & \textbf{Moody's} & \textbf{S\&P} \\
	\hline
	
	Focus
	&
	Indian securitisation market
	&
	Global default probability
	&
	Global structured finance
	\\
	\hline
	
	Primary Model
	&
	Portfolio Analysis
	&
	Idealised Default Rate
	&
	LEVELS Framework
	\\
	\hline
	
	Risk Assessment
	&
	Credit Enhancement
	&
	Expected Loss
	&
	Cash Flow + Stress Testing
	\\
	\hline
	
	Output
	&
	Credit Rating
	&
	Credit Rating
	&
	Credit Rating
	\\
	\hline
	
\end{tabularx}

\vspace{0.6cm}

\subsection*{Importance of Credit Ratings}

\begin{itemize}
	
	\item Help investors understand investment risk.
	
	\item Improve transparency in securitisation transactions.
	
	\item Support regulatory compliance.
	
	\item Determine pricing of securitised securities.
	
	\item Increase investor confidence.
	
\end{itemize}